% #2 MCFG v72 日本語査読用付録
% 数式・定理 label・証明の依存関係は英語版 appendix_window_abelian.tex と対応させる．
\section{固定窓と有限アーベル群を組み合わせた型付け}
\label{app:window-abelian}

この付録では，固定接頭辞--接尾辞窓と有限アーベル群の情報を組み合わせた場合に，
特徴データの総文字数についての多項式定理を拡張する．
群を用いない元の定理~\ref{thm:boundary-polydata} は変更しない．

\begin{lemma}[アーベル群値を指定した短いタプル]
\label{lem:group-typed-shortest}
$G=(V,\Sigma,P,S,\fan)$ を，ファンアウトが高々 $f$，
規則ランクが高々 $r_{\max}\ge1$ の縮約済み良形 MCFG とする．
$g:\Sigma^*\to A$ を，固定有限アーベル群への準同型とする．
各 $B\in V$ と，実現可能な任意のベクトル
$\mathbf s\in A^{\fan(B)}$ に対し，
\[
 \tuple u\in L_B(G),\qquad
 g^{(\fan(B))}(\tuple u)=\mathbf s
\]
を満たすタプルで，
\[
 \|\tuple u\|_1\le\beta_g:=
 (r_{\max}+1)\bar\theta_G
 \max\{1,(|A|^f-1)|V||A|^f\}
\]
となるものが存在する．
\end{lemma}

\begin{proof}
$B$ からの終端導出で，根の $g$-値ベクトルが $\mathbf s$ となるものを考える．
その中から，まず根の生成タプルの総文字数を最小化し，
次に導出木の節点数を最小化した導出 $T$ を選ぶ．
各節点に，その非終端記号 $C$ と，その部分木が生成するタプルの
成分ごとの $g$-値ベクトルを標識する．
この標識は高々 $n_g:=|V||A|^f$ 種類しかない．

ある祖先とその真の子孫が同じ標識をもつと仮定する．
祖先の部分木を子孫の部分木で置き換えると，
根の $g$-値ベクトルが保存される．
実際，MCFG の各規則は，子の成分別群値から親の群値を計算する．
非消去性により生成タプルの総文字数は増えず，
導出木の節点数は真に減る．
これは $T$ の選び方に矛盾する．
したがって木の深さは高々 $n_g$ である．

部分木が生成するタプルの総文字数が
$\bar\theta_G$ を超える節点を\emph{重い節点}と呼ぶ．
非消去性により，重い節点の祖先もすべて重い．
互いに比較不能な重い節点を
$v_1,\ldots,v_m$ とし，根のファンアウトを $d\le f$ とする．
各 $v_i$ の標識を $C_i$ とする．
補題~\ref{lem:thickness-local-bounds}(ii) により，
$C_i$ から長さ高々 $\bar\theta_G$ のタプルを生成する導出が存在する．
$v_i$ の部分木をこの短い導出で置き換えれば，
根の総文字数は真に減る．

文法は線形かつ非消去であるから，互いに比較不能な部分木の出力成分は，
根の各成分内で互いに交わらない終端ブロックに寄与する．
ただし成分ごとにその出現順序が異なることはある．
アーベル群では積を加法的に表せ，ブロックの順序に依存しない．
したがって，各節点の置換による変化を
$\delta_i\in A^d$ とおくと，
任意の部分集合 $J$ に対応する同時置換は
根の $g$-値ベクトルを $\sum_{i\in J}\delta_i$ だけ変化させる．
$m\ge|A|^d$ なら，$m+1$ 個の部分和
\[
 0,\quad\delta_1,\quad
 \delta_1+\delta_2,\quad\ldots,\quad
 \delta_1+\cdots+\delta_m
\]
のうち二つが一致する．
よって非空の連続部分集合に対応する置換は根の群値を保つ一方，
生成タプルを真に短くする．これは最小性に矛盾する．
したがって重い節点からなる任意の反鎖の大きさは
$|A|^d\le|A|^f$ 未満である．

根が軽ければ生成タプルの総文字数は
$\bar\theta_G$ 以下である．
根が重い場合，重い節点たちは祖先について閉じた木をなす．
その葉は重い節点の反鎖であり，各重い節点はこのいずれかの葉への経路上にある．
各経路の長さは高々 $n_g$ なので，
重い節点は高々 $(|A|^f-1)n_g$ 個である．
補題~\ref{lem:thickness-local-bounds}(i) により，
各規則が明示的に書く終端因子の総長は
$\bar\theta_G$ 以下である．
重い節点の軽い子はそれぞれ高々
$\bar\theta_G$ 個の終端記号に寄与し，
重い節点の子の個数は高々 $r_{\max}$ である．
重い節点とその軽い子を数えれば所望の上界が得られる．
$|A|=1$ の場合は，根が軽い場合の議論で扱える．
\end{proof}

\begin{lemma}[窓・アーベル群型付けの短い正準タプル]
\label{lem:window-abelian-short}
補題~\ref{lem:group-typed-shortest} の仮定の下で，
$k,\ell\ge0$ を固定する．
$h:\Sigma^*\to M$ を
$h\preceq(h_\Sigma^{k,\ell},g)$
を満たす有限型付けとする．
$\widetilde G_0$ を $G$ の完全 $h$-型付き精密化を刈り込んだものとし，
その開始記号以外の非終端記号数を $n_{\mathrm{typ}}$ とする．
このとき各非終端記号 $X$ について，正準タプル
$\omega(X)$ を
\[
 \|\omega(X)\|_1\le
 f(k+\ell+1)n_{\mathrm{typ}}\,
 \bigl(\bar\theta_G+r_{\max}\beta_g\bigr)
\]
となるように選べる．
\end{lemma}

\begin{proof}
$X$ からの型付き終端導出のうち，節点数が最小のものを選ぶ．
一つの経路上で型付き非終端記号が反復すれば，
根の $h$-型を保ったまま反復部分を除けるので，
各経路の節点数は高々 $n_{\mathrm{typ}}$ である．
補題~\ref{lem:boundary-short-canonical} の証明と同様に，
根の各成分で $h_\Sigma^{k,\ell}$ の型を決める位置を標識する．
成分が窓の閾値より短い場合は全位置を標識し，
それ以外は先頭 $k$ 位置と末尾 $\ell$ 位置を標識する．
$k=\ell=0$ の場合は1位置を標識する．
これらの標識への祖先経路がつくる骨格の節点数は高々
$f(k+\ell+1)n_{\mathrm{typ}}$ である．

骨格に含めなかった子部分木は，
それが寄与する根の各成分において，
標識されていない内部部分に完全に収まる．
その非終端記号 $C$ と，実際の成分別群値
$\mathbf s$ を記録する．
補題~\ref{lem:group-typed-shortest} により，
$C$ から同じ $\mathbf s$ を生成し，
総文字数が $\beta_g$ 以下の導出で置き換える．
良形性により置換後の全成分は非空に保たれる．
保護された境界位置は変化せず，
もとの長い成分は置換後も長いままである．
よって窓型は保存される．
また，根の各成分は規則中の明示的な終端因子と子成分の連接からなる．
置換した子成分の群値はそれぞれ保存されているので，
根の $g$-値ベクトルも変化しない．
この段階では群の可換性を新たに使う必要はない．

したがって積型 $(h_\Sigma^{k,\ell},g)$ の値が保存され，
それを通じて決定されるより粗い $h$-型も保存される．
中間の型付き標識はいったん忘れ，
得られた $G$ 上の導出を
命題~\ref{prop:output-invariants}(ii) で
$h$-型付き精密化に一意に持ち上げる．
$X$ は刈り込み後に残る非終端記号であるため，
この根をもつ導出も刈り込み後の文法に残る．
骨格の各節点が明示的に書く終端記号は
高々 $\bar\theta_G$ 個であり，
骨格から外れた子は高々 $r_{\max}$ 個である．
これらを数えて表示した長さ上界を得る．
\end{proof}

\begin{theorem}[窓・アーベル群型付けに対する多項式特徴データ]
\label{thm:window-abelian-polydata}
$f,r_{\max},k,\ell,\Sigma$，
有限アーベル群 $A$ と準同型 $g:\Sigma^*\to A$，
および有限型付け $h\preceq(h_\Sigma^{k,\ell},g)$ を固定する．
ファンアウトが高々 $f$，規則ランクが高々 $r_{\max}$ の
縮約済み良形文法 $G$ が
$\Cmcf{f}{h}$ に属する言語を生成するとする．
このとき，ランクを制限した再構成作用素
$\mathcal B^{\le r_{\max}}_{f,h}$ は，
次を満たす特徴標本をもつ．
\[
 |\Wit(\widetilde G_0)|=O(\|G\|),\qquad
 \|\Wit(\widetilde G_0)\|_+
   =O_{f,r_{\max},k,\ell,h,g,\Sigma}
     (\|G\|^3\bar\theta_G).
\]
\end{theorem}

\begin{proof}
固定有限型付けにより
$n_{\mathrm{typ}}\le|V||M|^f=O(\|G\|)$ である．
補題~\ref{lem:group-typed-shortest} から
$\beta_g=O(\|G\|\bar\theta_G)$ が従う．
補題~\ref{lem:window-abelian-short} により，
選択した正準タプルの総文字数は
$O(\|G\|^2\bar\theta_G)$ 以下である．
補題~\ref{lem:untyped-contexts} を適用すると，
各特徴証拠語の長さも
$O(\|G\|^2\bar\theta_G)$ 以下となる．
命題~\ref{prop:witness-cardinality} から
その語数は $O(\|G\|)$ である．
空語 $\eps$ の証拠が必要でも，
符号化された総長の増加は高々定数である．
\end{proof}

\begin{remark}
可換性を使うのは，
補題~\ref{lem:group-typed-shortest} の反鎖に対する零和の段階である．
ファンアウト1では各部分木は一つの区間に寄与し，
関連する CFG 論文は任意の有限群を扱っている．
一方，$f\ge2$ で群が非可換な場合，
異なる部分木の寄与が根の各成分で異なる順序に交錯し得る．
本稿の証明はこの場合を扱わないが，
不可能性を主張するものではない．
\end{remark}
